\section{A Unified Theory of Sleep Function and Consciousness}\label{sec:unified_theory}

The convergence of this evidence allows for a unified, multi-stage model of sleep's function. The different sleep stages are not redundant but are sequential, complementary components of a single nightly learning algorithm:
\begin{enumerate}
    \item \textbf{Wake:} Data Acquisition.
    \item \textbf{Hypnagogia (N1):} Procedural Skill Rehearsal (from Sec 2).
    \item \textbf{NREM (SWS):} Episodic Memory Consolidation \& Replay.
    \item \textbf{REM:} Generalization \& Semantic Extraction (via adversarial simulation).
\end{enumerate}

This biological imperative is directly linked to the nature of consciousness itself. Modern theories, such as the Global Neuronal Workspace (GNW) theory, model consciousness as the product of a "rapidly changing neural network" defined by "recurrent activity" \citep{dehaene1998gnw}. Sleep, therefore, can be understood as the essential, universal, offline maintenance and training algorithm that prunes, consolidates, and generalizes the data within the very neural network that, when "online," generates conscious experience.
